\section{Schema/Interface Analysis}

The current artifact uses the clarified schema for all main results. The earlier cued construction motivated a schema-sensitivity question, but its model results are quarantined and are not used here. After the main role-uncued pilot, we ran a small Phase 1.1 ablation over generated-lore and buried-primary rows. The ablation holds the evidence packet fixed and varies only the structured output interface. It is exploratory follow-up evidence, not a replacement for the Phase 1 main table.

The interface finding that remains valid in this pilot is narrower. Structured fields matter because evidence and action fields can fail while the final verdict succeeds. In generated-lore rows, strict failures are highly sensitive to whether the scorer treats polluted documents in support fields as accepted support or recognizes that prose and actions diagnosed them as unverified. In active-verification rows, the action schema makes executable next steps measurable instead of leaving them as prose.

The core-claim overnight rerun gives a second interface-sensitivity warning. A role-decomposed schema over 144 generated-lore rows changed operational contract success only slightly on average (mean delta 0.014) and did not create a new model ranking. It did, however, make the contract easier to interpret by separating clean support, refuting evidence, rejected or contaminated evidence, diagnostic evidence, and verification actions. This should be read narrowly: schema wording and field semantics are experimental factors for machine-readable evidence roles. They do not solve all strict-score artifacts, explain every Phase 1 gap, or use the quarantined cued results.
